\chapter*{ТЕРМИНЫ И ОПРЕДЕЛЕНИЯ}
\addcontentsline{toc}{chapter}{ТЕРМИНЫ И ОПРЕДЕЛЕНИЯ}
\setlength{\parindent}{0pt}
% Перечень в алфавитном порядке, без знаков препинания в конце строки (п. 4.9.2)
% Термины пишутся без жирного шрифта
Большая языковая модель --- нейросетевая модель на основе трансформерной архитектуры, обученная на больших текстовых корпусах и способная генерировать связный текст\\[0.5em]
Задача 2-4-6 --- психологическая парадигма для исследования индуктивного открытия правил, в которой испытуемый угадывает скрытое правило, предлагая числовые тройки и получая бинарную обратную связь\\[0.5em]
Задача выбора карточек Уэйсона --- психологическая парадигма для исследования условного рассуждения и ошибок логического вывода, в которой испытуемый выбирает карточки, необходимые для проверки условного правила\\[0.5em]
Когнитивное предубеждение --- устойчивое систематическое отклонение от нормативных стандартов вероятностного или логического рассуждения\\[0.5em]
Контаминация обучающих данных --- явление, при котором тестовые задачи присутствуют в обучающем корпусе модели, что искажает оценку её реальных когнитивных способностей\\[0.5em]
Ошибка конъюнкции --- когнитивное предубеждение, при котором вероятность конъюнкции двух событий оценивается как более высокая, чем вероятность одного из них\\[0.5em]
Предубеждение подтверждения --- когнитивное предубеждение, при котором агент предпочитает искать и принимать информацию, подтверждающую имеющуюся гипотезу, а не опровергающую её\\[0.5em]
Эвристика репрезентативности --- когнитивная стратегия оценки вероятности события по степени его соответствия типичному прототипу без учёта базовых вероятностей